\documentclass[12pt]{article}
\usepackage[utf8]{inputenc}
\usepackage[T1]{fontenc}
\usepackage{hyperref}
\title{Wergaia Language Resources for K-12 Curriculum Development}
\author{Compiled by Claude Sonnet 5 (claude-sonnet-5) for Andr\'as Kornai}
\date{}
\begin{document}
\maketitle
\section{Introduction}
Wergaia (ISO 639-3 weg) is an Aboriginal language of the Wimmera region of north-western Victoria, Australia. The Wotjobaluk community is running a revival effort. It is related to, and overlaps with, the Wotjobaluk dialect (xwt).
\section{Primary Educational Resources}
The VCE subject \emph{Indigenous Languages of Victoria: Revival and Reclamation} has reference materials that include a Wergaia selection, and the Victorian School of Languages has supported Wergaia learners. I saw no direct URL for either.
\section{Linguistic and Reference Materials}
\begin{itemize}
\item Wergaia language, Wikipedia: \url{https://en.wikipedia.org/wiki/Wergaia_language}
\item Julie Reid, ``The rebirth of Wergaia: a collaborative effort'': \url{https://open.sydneyuniversitypress.com.au/9781920899554/rl-the-rebirth-of-wergaia-a-collaborative-effort.html}
\item Increasing the accessibility of information on the Indigenous languages of Victoria: \url{https://open.sydneyuniversitypress.com.au/9781920899554/rl-increasing-the-accessibility-of-information-on-the-indigenous-languages-of-victoria.html}
\item Victorian Collections entries for Wergaia, including the Reid \emph{Wergaia Community Grammar and Dictionary} (2007 consultation draft) and archival recordings held by VACL: \url{https://victoriancollections.net.au/collections?q=wergaia}
\end{itemize}
\section{Religious and Community Resources}
The Wotjobaluk Knowledge Place in Dimboola is the community base for the revival. I saw no website for it.
\section{Digital and Audio Resources}
VACL holds 1960s recordings of Wergaia with Walter Kennedy and Billy Marks, per the Victorian Collections entries above.
\section{Curriculum Development Context}
Few speakers remain, so the curriculum would be revival-based. The community must lead, and access to the archival materials should be arranged with VACL and the Wotjobaluk community.
\end{document}